\chapter{What Kind of Problem Is This?}
\label{ch:problem}

Every modeling decision you will make --- how much data to train on, which
architecture to reach for, whether a $+0.0003$ improvement is real ---
depends on facts about the \emph{problem class} that are knowable before you
fit a single model. This chapter establishes those facts for the tabular
time-series genre, using the Jane Street competition as the running specimen.

\section{Anatomy of a modern market-forecasting competition}

The Jane Street \emph{Real-Time Market Data Forecasting} challenge (2024--25)
is a nearly perfect exemplar of the genre, so it is worth dissecting its
structure piece by piece. You will meet the same skeleton, with different
skin, in DRW's crypto challenge, Optiver's contests, and most proprietary
research problems inside quantitative trading firms.

\subsection{The panel}

The data is a \emph{panel}: observations indexed by three coordinates,
\[
(\text{\code{symbol\_id}},\ \text{\code{date\_id}},\ \text{\code{time\_id}})
\;\in\; \{0,\dots,38\} \times \{0,\dots,1698\} \times \{0,\dots,967\}.
\]
Roughly forty instruments, about seventeen hundred trading days, and 968
intra-day timestamps per day (after day 677, when the day length stabilizes;
before that the panel is ragged --- a detail that will matter in
Chapter~\ref{ch:preprocessing}). The full table is about 47 million rows.

Panels are richer than single time series in a specific, exploitable way:
at every timestamp you observe a \emph{cross-section} of symbols. Anything
common to the cross-section (``the market moved'') can be separated from
what is idiosyncratic to a symbol. Chapter~\ref{ch:featurepatterns} builds
features from exactly this decomposition; here it is enough to note that the
panel structure is not incidental --- it is one of the two or three main
places where signal hides.

\subsection{The features}

Each row carries 79 anonymized features, \feat{00} through \feat{78}, plus a
per-row importance weight. Anonymization is standard practice --- the host
cannot reveal what the columns mean --- and it sets up the central epistemic
game of the genre: \textbf{the data-generating process is hidden but not
unknowable}. The features have autocorrelation structures, cross-sectional
correlation patterns, tag metadata, and timing relationships with the target,
all of which survive anonymization. Part~II of this book is about
interrogating them.

A metadata file, \code{features.csv}, assigns each feature a subset of
seventeen boolean tags. Tags are the host whispering hints: features sharing
\code{tag\_12} tend to be transformations of one underlying quantity.
We will use tag structure repeatedly.

\subsection{The targets}

Nine anonymized \emph{responders}, \resp{0} through \resp{8}, of which
\resp{6} is the competition target. The others are given during training but
\emph{not} at inference time --- except through a lag mechanism described
below. Keep this asymmetry in mind: auxiliary series that exist in training
but not at test time are not useless; they are \emph{free auxiliary
supervision} (Chapter~\ref{ch:featurepatterns}) and, lagged, they are
features.

\subsection{The protocol}

This competition scores you through a \emph{streaming gateway}: your model
receives the feature rows for one \code{time\_id} (all symbols at once), must
return predictions before seeing the next timestamp, and --- crucially ---
receives the \emph{previous day's complete responder table} at the start of
each new day (the ``lags'' file). The protocol encodes the real-world
information set of a live trading system:

\begin{itemize}
\item Within a day, you never observe your target. No within-day
  autoregression on the label is possible.
\item Across days, you observe the target with a one-day delay, at full
  intra-day resolution. Yesterday's realized path is the freshest label
  information you will ever have.
\item Features arrive in real time and may be used contemporaneously.
\end{itemize}

We call this the \emph{deployment information set}, and
Chapter~\ref{ch:leakage} elevates ``respect the deployment information set''
into the central commandment of the genre. Architecturally, it is also why
models that treat \emph{yesterday's target path} and \emph{today's streaming
features} as two different kinds of input (Chapter~\ref{ch:attention}, the
TimeXer design) map so naturally onto this problem.

\section{The noise ceiling: the number that rules everything}
\label{sec:noiseceiling}

Here is the most important table in this book. It is the final private
leaderboard of the Jane Street competition --- the scores achieved on months
of genuinely unseen future data:

\begin{center}
\begin{tabular}{@{}rlr@{}}
\toprule
place & team & private weighted $R^2$ \\
\midrule
1 & ms capital & $0.01389$ \\
2 & Patrick Yam & $0.01327$ \\
3 & (undisclosed) & $0.01316$ \\
4 & Haoze Hou & $0.01168$ \\
8 & Evgeniia Grigoreva & $0.01043$ \\
19 & John Payne & $0.00930$ \\
49 & --- & $0.00813$ \\
\bottomrule
\end{tabular}
\end{center}

Read it slowly. The best result achieved by roughly 2{,}800 competing teams,
many of them professionals with serious compute, was to explain
\textbf{1.39 percent} of the target's variance. The difference between a gold
medal and 49th place was about half a percentage point of $R^2$. And --- as
we will derive in Chapter~\ref{ch:targets} --- this is not because everyone
was bad at the problem. The target is constructed (a forward moving average
of near-white innovations) so that most of its variance is \emph{physically
unpredictable} from the available information. There is a ceiling, the
ceiling is low, and everybody is compressed underneath it.

\begin{keybox}[: estimate the ceiling first]
Before modeling, estimate how much of the target is knowable at all. Use the
leaderboard of similar past competitions, the autocorrelation structure of
the target (Chapter~\ref{ch:targets}), or an explicit synthetic replica of
the data-generating process (Chapter~\ref{ch:synthetic}). Every downstream
judgment --- what improvement is plausible, what improvement is
\emph{suspicious} --- calibrates against this number.
\end{keybox}

\subsection{Consequence 1: variance reduction beats capacity}

Suppose the predictable component of the target is a function $f$ with
$\Var(f) \approx 0.014 \cdot \Var(y)$, and your model produces
$\hat f = f + \varepsilon$ where $\varepsilon$ is estimation error ---
error from finite data, noisy optimization, unlucky seeds. Your achieved
score is roughly
\[
R^2 \;\approx\; \frac{\Var(f) - \E\|\varepsilon\|^2_{\text{aligned}}}{\Var(y)},
\]
that is, \emph{estimation variance subtracts directly from a signal budget
that was tiny to begin with}. When the budget is 60\% of variance (image
classification), estimation noise is a rounding error. When the budget is
1.4\%, estimation noise is often \emph{larger than the signal}, and the
highest-return activities are the ones that shrink $\varepsilon$: averaging
over seeds, bagging over time windows, ensembling decorrelated models,
online refitting toward the current regime.

Our campaign measured this directly. The best single model (an LSTM with
auxiliary heads, Chapter~\ref{ch:rnn}) scored $+0.0092$; averaging it with
two other models --- one of them a gradient-boosted tree scoring only
$+0.0068$ --- produced $+0.0104$. Meanwhile, doubling the LSTM's hidden
width, a capacity move that would matter enormously in a high-SNR domain,
changed the score by less than seed noise. \emph{At low SNR, the ensemble is
the model.}

\subsection{Consequence 2: your validation score is a random variable}

The same arithmetic applies to evaluation. A validation set of 100 trading
days ($\sim$3.7 million rows, but with heavy temporal correlation --- the
effective sample size is closer to the number of days than the number of
rows) measures a true score of $\sim 0.01$ with sampling noise on the order
of a few $10^{-4}$. Two consequences:

\begin{itemize}
\item A measured difference of $+0.0002$ between two configurations is,
  by itself, \emph{nothing}. Demand replication: multiple seeds, multiple
  validation windows, or both.
\item Model selection over many candidates on one validation window is
  itself an overfitting process. If you compare thirty ideas on the same
  100-day tail, the winner's margin is partly selection bias. Hold back a
  second window for the finalists (Chapter~\ref{ch:validation}).
\end{itemize}

\begin{jsbox}[: seed noise vs.\ claimed effects]
In our runs, retraining the same GRU with a different random seed moved the
validation weighted $R^2$ by $\pm 0.0003$ to $\pm 0.0008$. Nearly every
``improvement'' under $0.0005$ that we chased down was seed noise wearing a
costume. The discipline that survived: any effect claimed at under
$0.001$ must show up in at least two seeds \emph{and} two validation
windows before it earns a line in the technical report.
\end{jsbox}

\subsection{Consequence 3: overfitting is the default outcome, not a risk}

A model with a few hundred thousand parameters, trained on a target whose
predictable share is 1.4\%, has abundant capacity to memorize noise patterns
that \emph{look} like that 1.4\%. In-sample $R^2$ of 0.05--0.30 with
out-of-sample $R^2$ of 0.005 is not a pathology; it is the equilibrium state
of this genre. All of Part~III --- leakage discipline, purged validation,
frozen preprocessing statistics --- exists because at these SNRs the standard
machine-learning failure modes are not occasional accidents but the
\emph{default trajectory} of any project that does not actively resist them.

\section{Drift: the second law of the genre}
\label{sec:drift}

Financial panels are not stationary. The relationship between features and
target decays over time --- regimes change, market participants adapt, the
host's data pipeline evolves. Our cleanest measurement of this
(Chapter~\ref{ch:online} has the full story): a GRU trained on 800 days of
history scored \emph{worse} on the recent validation tail ($+0.0077$) than
the same architecture trained on the most recent 280 days ($+0.0089$).
More data made the model worse, because the older data taught relationships
that no longer hold. The static score of old-data models on recent days is
frequently \emph{negative} --- worse than predicting zero.

Drift has two practical faces:

\begin{enumerate}
\item \textbf{Recency windowing.} There is an optimal training window ---
  long enough to estimate, short enough to be current. Sweep it
  (Chapter~\ref{ch:online}); in our problem the knee was around 300--400
  days.
\item \textbf{Online adaptation.} If the protocol delivers labels with a
  delay (here: one day), a model that takes a small gradient step on each
  day's newly revealed labels tracks the drifting relationship. This single
  mechanism was worth more than any architectural choice we made:
  $+0.003$ to $+0.005$ weighted $R^2$, on every recurrent model, in every
  configuration we tested.
\end{enumerate}

\section{The mindset}

Let us compress the chapter into an operating posture. You are entering a
domain where:

\begin{itemize}
\item the winner explains 1.4\% of variance, so \textbf{humility is
  quantitative}: any result that looks much better than the ceiling is a bug
  (almost always leakage, Chapter~\ref{ch:leakage});
\item measurement noise rivals effect sizes, so \textbf{replication is not
  optional};
\item the world drifts, so \textbf{recency and adaptation beat volume};
\item signal is scarce but structure is everywhere, so \textbf{forensics
  (Part II) pays better than architecture search (Part V)};
\item and estimation variance is the dominant error term, so
  \textbf{boring averaging wins} (Chapter~\ref{ch:ensembling}).
\end{itemize}

The rest of the book is the long version of this list.

\begin{drillbox}
Before reading further, take any time-series competition you have access to
and, without fitting a model: (1) find the metric's zero-point --- what score
does predicting all zeros achieve? (2) find the winner's score and express it
as a fraction of the theoretical maximum; (3) compute the target's lag-1
autocorrelation. These three numbers place the competition in the genre's
coordinate system and predict most of what will and won't work.
\end{drillbox}
